\documentclass[11pt]{article}
\usepackage[utf8]{inputenc}
\usepackage{amsmath,amsthm,amsfonts,amssymb,amscd}
\usepackage{multirow,booktabs}
\usepackage[table,dvipsnames]{xcolor}
\usepackage{fullpage}
\usepackage{lastpage}
\usepackage{enumitem}
\usepackage{fancyhdr}
\usepackage{mathrsfs}
\usepackage{wrapfig}
\usepackage{listings}

\usepackage[round,authoryear]{natbib}

\usepackage{setspace}
\usepackage{calc}
\usepackage{multicol}
\usepackage{cancel}
\usepackage[retainorgcmds]{IEEEtrantools}
\usepackage[margin=3cm]{geometry}
\usepackage{amsmath}
\newlength{\tabcont}
\setlength{\parindent}{0.0in}
\setlength{\parskip}{0.05in}
\usepackage{empheq}
\usepackage{framed}
\usepackage[most]{tcolorbox}
\usepackage{xcolor}
\colorlet{shadecolor}{orange!15}
\parindent 0in
\parskip 1pt
\geometry{margin=1in, headsep=0.25in}
\theoremstyle{definition}
\newtheorem{defn}{Definition}
\newtheorem{reg}{Rule}
\newtheorem{exer}{Exercise}
\newtheorem{note}{Note}
\usepackage{tikz}

\definecolor{MyIndigo}{HTML}{DC143C}
\definecolor{linkcolor}{HTML}{DC143C}

\usepackage[
    colorlinks = true,
    linkcolor  = linkcolor,
    citecolor  = linkcolor,
    urlcolor   = linkcolor,
]{hyperref}

\setlength{\parindent}{0pt}
\setlength{\parskip}{8pt plus 2pt minus 1pt}

\usetikzlibrary{decorations.text}

\usepackage[T1]{fontenc}
\usepackage[sc]{mathpazo}
\linespread{1.05}
\usepackage{microtype}

\newcommand{\totalpages}{\pageref{LastPage}}
\newtheorem{remark}{Note}

\definecolor{codegreen}{rgb}{0,0.6,0}
\definecolor{codegray}{rgb}{0.5,0.5,0.5}
\definecolor{codepurple}{rgb}{0.58,0,0.82}
\definecolor{backcolour}{rgb}{0.95,0.95,0.92}

\newcommand{\s}{\sum_{i=1}^n}
\usepackage{bm}
\usepackage{tabularx}
\newcommand{\qq}{\quad\quad}

\usepackage[rightcaption]{sidecap}
\usepackage{graphicx}
\hypersetup{
    colorlinks = true,
    linkcolor  = MyIndigo,
    citecolor  = MyIndigo,
    urlcolor   = MyIndigo
}

\definecolor{MyIndigo}{HTML}{2F3A8F}

\begin{document}

\thispagestyle{empty}

\begin{minipage}[c]{0.48\textwidth}
  \centering
  {\LARGE\bf Datacenter LP, \textit{notes}}\\
  \vspace{2mm}
  {\large \textit{KK vs.\ VE},}
  Linus Lindquist
\end{minipage}
\hfill
\begin{minipage}[c]{0.48\textwidth}
  \textit{\textbf{Goal.} Document the LP model used to find the least-cost
  solar--wind--battery mix for a 1\,GW datacenter subject to an hourly
  clean-energy constraint. Covers variables, objective, constraints, economic
  intuition, and the solver.}
\end{minipage}

\rule{\textwidth}{0.8pt}

% ─────────────────────────────────────────────────────────────────────────────
\subsection*{Problem setting}

A datacenter draws a \textit{flat} $P = 1{,}000$\,MW around the clock, every
hour of the year ($T = 8{,}760$ hours). Danish regulation requires that a
fraction $x \in (0,1)$ of this demand be sourced \textit{on-site} in each
individual hour; the remainder may be imported from the Nordic spot market at an
exogenous hourly price $p_t$ (DKK/MWh, converted to EUR at $7.46$). This
hourly on-site requirement is referred to throughout as the
\textbf{Carbon-Free Energy (CFE)} constraint: in every hour $t$, grid purchases
$g^+_t$ may not exceed the grid cap $\bar{g} = (1-x)\cdot P$.

The model compares two on-site technologies on \textbf{total annualised cost}:
\begin{itemize}[noitemsep]
  \item \textbf{KK} — a 1\,GW SMR running flat with a single planned outage of
        $10\%$ of the year.
  \item \textbf{VE} — jointly optimised solar PV, onshore wind, and a
        lithium-ion BESS.
\end{itemize}

This document focuses on the VE case, which is formulated as a single linear
programme solved by HiGHS.

% ─────────────────────────────────────────────────────────────────────────────
\subsection*{Assumptions}

All cost and technical parameters are sourced from the Danish Energy Agency
Technology Data (DEA, 2030-projection column) and internal ET estimates.
Costs are in EUR; DKK figures converted at $7.46$\,DKK/€.
The capital recovery factor is $\text{CRF}(r,n) = r(1+r)^n/[(1+r)^n-1]$
at $r = 4\%$.

\medskip

{\small
\begin{tabularx}{\linewidth}{l >{\raggedright\arraybackslash}X r l l}
\toprule
Technology & Parameter & Value & Unit & Source \\
\midrule
\multirow{3}{*}{Common}
  & Discount rate $r$       &   4.00 & \%         & ET \\
  & Exchange rate           &   7.46 & DKK/€      & ET \\
  & Demand $P$              & 1{,}000 & MW        & — \\
\midrule
\multirow{6}{*}{SMR (KK)}
  & Capex                   & 8{,}000 & k€/MW     & ET \\
  & Fixed O\&M              &       0 & €/MW/yr   & ET \\
  & Variable O\&M (total)   &   31.08 & €/MWh     & ET \\
  & \quad of which: fuel    &    9.32 & €/MWh     & ET \\
  & \quad of which: O\&M    &   20.42 & €/MWh     & ET \\
  & \quad of which: surcharge&   1.34 & €/MWh     & ET \\
  & Lifetime                &      60 & yr        & FOA \\
  & Planned downtime $\delta$&     10 & \%        & ET \\
  & CRF$(4\%,\,60)$         &    4.42 & \%        & derived \\
\midrule
\multirow{5}{*}{Solar PV}
  & Capex                   &     450 & k€/MW     & DEA 2030 \\
  & Fixed O\&M (total)      &  15{,}200 & €/MW/yr & DEA 2030 + ET \\
  & \quad of which: DEA O\&M&  10{,}400 & €/MW/yr & DEA 2030 \\
  & \quad of which: land rent&  4{,}800 & €/MW/yr & ET \\
  & Variable O\&M           &       0 & €/MWh     & DEA 2030 \\
  & Lifetime                &      40 & yr        & DEA 2030 \\
  & CRF$(4\%,\,40)$         &    5.05 & \%        & derived \\
\midrule
\multirow{5}{*}{Wind (onshore)}
  & Capex                   &  1{,}150 & k€/MW    & DEA 2030 \\
  & Fixed O\&M              &  16{,}663 & €/MW/yr & DEA 2030 \\
  & Variable O\&M $\eta$    &    1.98 & €/MWh     & DEA 2030 \\
  & Lifetime                &      30 & yr        & DEA 2030 \\
  & CRF$(4\%,\,30)$         &    5.78 & \%        & derived \\
\midrule
\multirow{6}{*}{Battery (BESS)}
  & Power capex             &      80 & k€/MW     & DEA 2030 \\
  & Energy capex            &     200 & k€/MWh    & DEA 2030 \\
  & Fixed O\&M (on power)   &   8{,}100 & €/MW/yr & DEA 2030 \\
  & Lifetime                &      20 & yr        & DEA 2030 \\
  & Duration sh             &       4 & h         & assumed \\
  & CRF$(4\%,\,20)$         &    7.36 & \%        & derived \\
\midrule
\multirow{4}{*}{Grid tariffs}
  & Buy tariff $\tau^+$     &   11.66 & €/MWh     & Energinet \\
  & \quad (= 8.70\,øre/kWh) &        &            & \\
  & Sell tariff $\tau^-$    &    1.54 & €/MWh     & Energinet \\
  & \quad (= 1.15\,øre/kWh) &        &            & \\
\midrule
\multirow{3}{*}{Grid connection}
  & Fixed component         &   18.20 & M\,DKK    & ET \\
  & Variable component      &    0.663 & M\,DKK/MW & ET \\
  & Annualisation           & perpetuity at $r=4\%$ && ET \\
\bottomrule
\end{tabularx}}

\medskip
\textit{Note on SMR variable O\&M:} the fuel cost ($9.32$\,€/MWh) already
accounts for thermal efficiency, i.e.\ it is expressed per MWh of electrical
output. The ET surcharge ($1.34$\,€/MWh $\approx 1$\,øre/kWh) is an
additional policy cost applied on top of DEA estimates.

% ─────────────────────────────────────────────────────────────────────────────
\subsection*{Notation}

{\small
\begin{tabularx}{\linewidth}{l l >{\raggedright\arraybackslash}X}
\toprule
Symbol & Unit & Meaning \\
\midrule
$c_s,\,c_w$ & MW & installed solar / wind capacity \\
$P_b,\,E_b$ & MW, MWh & battery power and energy capacity \\
$u_t,\,d_t$ & MW & battery charge / discharge at hour $t$ \\
$s_t$ & MWh & battery state of charge (end of hour $t$) \\
$g^+_t,\,g^-_t$ & MW & grid buy / sell at hour $t$ \\
$q^s_t,\,q^w_t$ & MW & solar / wind curtailment at hour $t$ \\
$e_t$ & MW & CFE excess slack (penalised) \\
$\bar{g}$ & MW & maximum grid import/export: $(1-x)\cdot P$ \\
$\text{CF}^s_t,\,\text{CF}^w_t$ & — & solar / wind capacity factor at hour $t$ \\
$p_t$ & €/MWh & spot price at hour $t$ \\
$\tau^+,\,\tau^-$ & €/MWh & grid buy / sell tariff ($11.66$ / $1.54$) \\
$\eta$ & €/MWh & wind variable O\&M ($1.98$) \\
$M$ & €/MWh & CFE penalty ($10^6$) \\
\bottomrule
\end{tabularx}}

% ─────────────────────────────────────────────────────────────────────────────
\subsection*{Decision variables}

The LP has $3 + 8T \approx 70{,}083$ decision variables (with $\text{sh} = 4$
fixed; setting $\text{sh} = \texttt{None}$ adds $E_b$ as a fourth free capacity
variable):

\medskip
{\small
\begin{tabularx}{\linewidth}{>{\raggedright\arraybackslash}X l}
\toprule
Block & Size \\
\midrule
Capacity: $c_s,\,c_w,\,P_b$ (with $E_b = 4P_b$) & 3 scalars \\
Battery charge $u_t$ & $T$ \\
Battery discharge $d_t$ & $T$ \\
Battery SOC $s_t$ & $T$ \\
Grid buy $g^+_t$ & $T$ \\
Grid sell $g^-_t$ & $T$ \\
Solar curtailment $q^s_t$ & $T$ \\
Wind curtailment $q^w_t$ & $T$ \\
CFE excess $e_t$ & $T$ \\
\bottomrule
\end{tabularx}}

\medskip
All variables are continuous and non-negative. The capacity variables are shared
across all hours — the LP simultaneously decides \textit{how much} to build and
\textit{how to dispatch} each asset in every hour of the year in a single solve.
This is the key advantage over a two-stage Nelder-Mead outer / LP inner
approach: no iteration is needed.

% ─────────────────────────────────────────────────────────────────────────────
\subsection*{Objective function}

The LP minimises the total annual cost:

\begin{align}
\min \quad &
  \underbrace{\alpha_s c_s + \alpha_w c_w}_{\text{solar and wind capex + O\&M}}
  +\; \underbrace{\beta_P P_b + \beta_E E_b}_{\text{battery capex + O\&M}}
  \notag \\[4pt]
&  + \sum_{t=1}^{T}
  \Bigl[
    \underbrace{(p_t + \tau^+)\,g^+_t}_{\text{grid purchase cost}}
    - \underbrace{(p_t - \tau^-)\,g^-_t}_{\text{grid revenue}}
    - \underbrace{\eta\,q^w_t}_{\text{wind O\&M saved}}
    + \underbrace{M\,e_t}_{\text{CFE penalty}}
  \Bigr]
\label{eq:obj}
\end{align}

\textbf{Capacity cost coefficients.} Each coefficient annualises capital and
fixed operating costs using the capital recovery factor
$\text{CRF}(r,n) = r(1+r)^n / [(1+r)^n - 1]$ at discount rate $r = 4\%$:

\begin{align*}
\alpha_s &= \text{capex}_s \cdot \text{CRF}_s + \text{opex}^f_s
            + \underbrace{0}_{\eta^s=0} \cdot \textstyle\sum_t \text{CF}^s_t \\
\alpha_w &= \text{capex}_w \cdot \text{CRF}_w + \text{opex}^f_w
            + \eta \cdot \textstyle\sum_t \text{CF}^w_t \\
\beta_P  &= \text{capex}^P_b \cdot \text{CRF}_b + \text{opex}^f_b \\
\beta_E  &= \text{capex}^E_b \cdot \text{CRF}_b
\end{align*}

Note that $\alpha_w$ charges wind variable O\&M on \textit{all potential}
generation ($c_w \cdot \sum_t \text{CF}^w_t$). The term $-\eta\, q^w_t$ then
credits back the O\&M that would have been paid on curtailed wind, so the net
effect is that only \textit{dispatched} wind MWh carry the 1.98\,€/MWh charge.
Solar has no variable O\&M ($\eta^s = 0$), so curtailing solar saves nothing
and $q^s_t$ has zero cost in the objective.

\textbf{Economic implication.} Because $q^w_t$ enters the objective with a
negative coefficient $-\eta < 0$, the LP has a genuine economic incentive to
curtail wind in preference to solar whenever it must dispose of surplus. It will
always curtail wind first — not as an ad-hoc rule, but as the optimum.

\textbf{Grid terms.} The LP buys from the grid at the spot price plus
$\tau^+ = 11.66$\,€/MWh (\textit{forbrugstarif}) and sells at the spot price
minus $\tau^- = 1.54$\,€/MWh (\textit{produktionstarif}). At negative spot
prices the sell revenue $p_t - \tau^-$ goes negative — but since grid\_sell is
non-negative and enters the objective with a ``minus'' sign, selling at negative
net revenue is never chosen unless forced by the energy balance.

% ─────────────────────────────────────────────────────────────────────────────
\subsection*{Constraints}

\subsubsection*{Energy balance (T equalities)}

Every hour the electricity in equals the electricity out:

\begin{equation}
  c_s\,\text{CF}^s_t + c_w\,\text{CF}^w_t
  + d_t - u_t
  + g^+_t - g^-_t
  - q^s_t - q^w_t = P
  \quad \forall\, t
  \label{eq:balance}
\end{equation}

Left-hand side: gross VE generation, net battery flow (positive = discharge),
net grid flow (positive = buy), minus curtailment. The right-hand side is the
fixed demand $P$. This must hold with equality — power cannot be ``stored'' in
the air; any surplus must be exported, stored, or curtailed.

\subsubsection*{SOC dynamics (T equalities)}

The battery state of charge evolves as:

\begin{equation}
  s_t = s_{t-1} + u_t - d_t
  \quad \forall\, t, \qquad s_0 = 0
  \label{eq:soc}
\end{equation}

Initialising at $s_0 = 0$ is a conservative worst-case assumption: the battery
is empty at the start of the year, so any deficit in the opening hours must be
covered by grid import or VE generation without battery help. In practice the
LP wraps around by choosing a dispatch profile that leaves the battery at a
useful level throughout.

\subsubsection*{CFE constraint (T inequalities)}

The legal requirement is that the datacenter sources at least fraction $x$ of
its electricity on-site \textit{each hour}. Equivalently, grid imports cannot
exceed the complement:

\begin{equation}
  g^+_t - e_t \leq \bar{g} = (1-x)\cdot P
  \quad \forall\, t
  \label{eq:cfe}
\end{equation}

The slack variable $e_t \geq 0$ absorbs CFE violations and is penalised at
$M = 10^6$\,€/MWh in the objective. Binding the constraint hard would risk LP
infeasibility in hours of extreme shortfall (e.g., dark winter nights with no
battery left); the soft formulation with a very large penalty always keeps the
LP feasible while ensuring $e_t \approx 0$ at the optimum. The residual
$\sum_t e_t$ is reported as \texttt{cfe\_shortfall\_mwh}.

\subsubsection*{Battery bounds (3T inequalities)}

\begin{equation}
  u_t \leq P_b, \qquad d_t \leq P_b, \qquad s_t \leq E_b
  \quad \forall\, t
  \label{eq:batt}
\end{equation}

Battery power and energy capacity are distinct physical components priced
separately — the inverter ($P_b$, MW) and the cell stack ($E_b$, MWh). Their
ratio is fixed by a configurable parameter: $E_b = \text{sh} \times P_b$
where $\text{sh} = 4$\,h is set exogenously, so only $P_b$ is a free LP
variable (reducing capacity variables to 3). Setting $\text{sh} = \texttt{None}$
promotes $E_b$ to a fourth free variable with its own capex term. The model
treats storage as lossless (no round-trip efficiency loss).

\subsubsection*{Curtailment bounds (2T inequalities)}

\begin{equation}
  q^s_t \leq c_s\,\text{CF}^s_t, \qquad q^w_t \leq c_w\,\text{CF}^w_t
  \quad \forall\, t
  \label{eq:curtail}
\end{equation}

You cannot curtail more of a source than it actually generates in that hour.
These are linear in the capacity variables because $\text{CF}^s_t$ and
$\text{CF}^w_t$ are exogenous parameters. Without these bounds the LP could, in
principle, assign unbounded curtailment credit to wind (since $-\eta < 0$);
these constraints prevent that.

\subsubsection*{Grid sell bound (T inequalities)}

\begin{equation}
  g^-_t \leq \bar{g}
  \quad \forall\, t
\end{equation}

Grid exports are capped at the same grid connection capacity as imports. This
prevents the model from building arbitrarily large VE for pure export revenue,
which would not be commercially realistic for an on-site facility.

\subsubsection*{Grid buy bound and LP feasibility}

The effective constraint is the CFE inequality \eqref{eq:cfe}: $g^+_t \leq
\bar{g}$ at the optimum (since $e_t = 0$ always holds when $M = 10^6$). The
\textit{box} bound in the implementation is, however, set to
\begin{equation}
  g^+_t \leq P \quad \forall\, t
\end{equation}
rather than $\bar{g}$. This is purely a numerical device to keep the LP
feasible during the simultaneous capacity-and-dispatch solve. Before capacities
are set, there may be no feasible dispatch with $g^+_t \leq \bar{g}$: a dark,
windless hour with an empty battery and zero installed VE has no feasible
solution at the hard cap. Allowing $g^+_t$ up to $P$ and penalising exceedance
via $e_t$ at $M = 10^6$\,€/MWh ensures the LP is always feasible. At any
economically rational optimum — where the penalty cost would dwarf all other
costs — $e_t = 0$ and the hard CFE constraint $g^+_t \leq \bar{g}$ is
\textit{never} violated. The residual $\sum_t e_t$ (reported as
\texttt{cfe\_shortfall\_mwh}) is zero to numerical precision.

% ─────────────────────────────────────────────────────────────────────────────
\subsection*{Economic intuition — what is the LP doing?}

Think of the LP as a \textbf{perfect-foresight central planner} who knows all
8{,}760 spot prices for the year and chooses, once, how much to build and how to
dispatch every hour.

\textbf{Capacity sizing.} The planner trades off capital cost against
operational flexibility. More solar and wind reduce grid purchases but require
more curtailment management and battery storage to satisfy the hourly CFE
constraint. More battery allows the planner to smooth VE variability and shift
cheap renewable energy to Dunkelflaute hours, but battery is expensive. The
optimal capacity mix is the one where the marginal cost of adding another MW of
any technology equals the marginal value of the dispatch flexibility it provides
across all 8{,}760 hours simultaneously.

\textbf{Dispatch logic in surplus hours.} When VE generation exceeds demand the
LP can: (i) charge the battery, (ii) export to the grid, or (iii) curtail.
Exporting earns revenue $(p_t - \tau^-)$ and is preferred when spot prices are
high. Charging the battery is free (no variable cost) and valuable when future
hours are expected to have a deficit. Curtailment is the last resort — but since
curtailing wind saves $\eta = 1.98$\,€/MWh in O\&M, wind is always curtailed
before solar.

\textbf{Dispatch logic in deficit hours.} When VE generation falls short of the
required on-site floor $x \cdot P$, the LP can: (i) discharge the battery or
(ii) accept a CFE penalty. Discharging costs nothing in variable terms; the
penalty $M$ is prohibitively large. In practice the LP sizes the battery large
enough to avoid all CFE violations.

\textbf{Interaction between curtailment and battery.} Because curtailing wind is
beneficial ($-\eta < 0$ in the objective), the LP has an incentive to use
battery capacity to absorb solar surplus (freeing up the energy balance to
curtail wind). This creates a feedback: larger battery $\Rightarrow$ more wind
curtailment enabled $\Rightarrow$ lower wind O\&M cost. The battery and
curtailment decisions are not separable.

% ─────────────────────────────────────────────────────────────────────────────
\subsection*{The LP solve — structure and solver}

\textbf{Matrix form.} The LP is written as

\[
  \min_{\mathbf{x}} \; \mathbf{c}^\top \mathbf{x}
  \quad \text{s.t.} \quad
  A_{\text{eq}}\, \mathbf{x} = \mathbf{b}_{\text{eq}}, \quad
  A_{\text{ub}}\, \mathbf{x} \leq \mathbf{b}_{\text{ub}}, \quad
  \mathbf{l} \leq \mathbf{x} \leq \mathbf{u}
\]

where all matrices are assembled in \texttt{scipy.sparse} CSR format. Sparsity
is critical: the energy balance and SOC constraints each touch only $O(1)$
non-zeros per row, giving a matrix $A_{\text{eq}}$ of size
$2T \times (4 + 8T)$ with roughly $8$ non-zeros per row rather than
$4 + 8T \approx 70{,}088$ in a dense representation.

The constraint blocks are:

\medskip
{\small
\begin{tabularx}{\linewidth}{>{\raggedright\arraybackslash}X l l}
\toprule
Block & Type & Rows \\
\midrule
Energy balance (\ref{eq:balance}) & equality & $T$ \\
SOC recurrence (\ref{eq:soc}) & equality & $T$ \\
CFE inequality (\ref{eq:cfe}) & inequality & $T$ \\
Battery charge bound (\ref{eq:batt}) & inequality & $T$ \\
Battery discharge bound (\ref{eq:batt}) & inequality & $T$ \\
Battery SOC bound (\ref{eq:batt}) & inequality & $T$ \\
Solar curtailment bound (\ref{eq:curtail}) & inequality & $T$ \\
Wind curtailment bound (\ref{eq:curtail}) & inequality & $T$ \\
\bottomrule
\end{tabularx}}

\medskip
Total: $2T$ equality rows and $6T$ inequality rows, plus simple box constraints
on each variable. With $T = 8{,}760$ this gives roughly $52{,}560$ constraints
over $70{,}084$ variables.

\textbf{Solver.} The problem is passed to \textbf{HiGHS} via
\texttt{scipy.optimize.linprog(..., method='highs')}. HiGHS is a
state-of-the-art open-source LP/MIP solver that uses a \textit{revised dual
simplex} method (with possible interior-point warm-start) and aggressive
sparsity exploitation. The solve typically completes in a few seconds on a
single thread.

\textbf{Perfect foresight.} The LP assumes the planner knows all spot prices
$p_1, \ldots, p_T$ at the time of dispatch. This is an idealisation: a real
operator would face uncertainty and solve a receding-horizon stochastic
programme. Perfect foresight gives a \textit{lower bound} on actual operating
costs, so VE costs reported here are optimistic relative to real-world
performance. KK does not face this issue — the SMR runs at constant output
and the outage window is chosen using the same perfect-foresight assumption
(minimise $\sum_{t \in \text{outage}} (p_t + \tau^+)$ over all contiguous
windows of length $0.10 \times T$).

% ─────────────────────────────────────────────────────────────────────────────
\subsection*{KK benchmark}

The SMR option requires no LP. Output is constant and the problem reduces to
placing a single maintenance window, making the cost fully analytical.

\textbf{Dispatch.} A single 1{,}000\,MW reactor runs at full output for every
hour of the year except a contiguous block of $\delta T = 0.10 \times 8{,}760 =
876$ hours of planned maintenance. Outside the outage, the reactor exactly
meets demand with no grid interaction; no battery or VE capacity is needed. The
SMR is sized at exactly $P = 1{,}000$\,MW — no oversizing.

\textbf{Outage placement.} The 876-hour window is placed at the cheapest
available hours to minimise grid expenditure during downtime. Define the total
grid cost for a window starting at hour $k$:
\[
  C(k) = \sum_{t=k}^{k+876-1} (p_t + \tau^+) \cdot P
\]
The optimal start is $k^* = \arg\min_k C(k)$, found in $O(T)$ via a cumulative
sum. In the 2025--26 data this places the outage in late May to late June, when
spot prices are low due to high Nordic hydro and solar output.

\textbf{CFE constraint during outage.} During the 876-hour outage the
datacenter must import $P = 1{,}000$\,MW — far above the hourly CFE cap
$\bar{g} = (1-x)\cdot P$. Planned maintenance is treated as a regulatory
exemption: the CFE constraint is waived for these hours. Grid connection costs
are internalised via the tilslutningsbidrag (see below), which is sized for the
full 1{,}000\,MW import.

\textbf{Cost components.} Total annual KK cost:
\begin{align*}
  C_{\text{KK}} &=
    \underbrace{\text{capex}_{\text{KK}} \cdot \text{CRF}_{\text{KK}} \cdot P
                + \text{opex}^v_{\text{KK}} \cdot P\,(1-\delta)T}_{\text{on-site: capex + fuel + O\&M}} \\
  &\quad
    + \underbrace{\sum_{t \in \mathcal{O}} (p_t + \tau^+) \cdot P}_{\text{grid purchase during outage}}
    + \underbrace{\text{gc}(P)}_{\text{grid connection}}
\end{align*}

where $\text{opex}^v_{\text{KK}} \approx 31.08$\,€/MWh covers fuel ($9.32$),
variable O\&M ($20.42$), and an ET policy surcharge ($1.34$), and
$\text{gc}(P)$ is the annualised tilslutningsbidrag for a $P = 1{,}000$\,MW
connection.

\textbf{No battery, no export.} KK never exports to the grid and carries no
battery. Outside the outage it is self-sufficient; during the outage it
relies entirely on grid import. This simplicity is also a vulnerability: the
outage cost is fully exposed to whatever spot prices prevail in the chosen
window, with no storage buffer.

% ─────────────────────────────────────────────────────────────────────────────
\subsection*{Outputs and comparison}

After solving, the LP returns:

\begin{itemize}[noitemsep]
  \item Optimal capacities: $c_s^*,\, c_w^*,\, P_b^*,\, E_b^*$.
  \item Full hourly dispatch arrays: $u_t,\, d_t,\, s_t,\, g^+_t,\, g^-_t,\,
        q^s_t,\, q^w_t,\, e_t$ for all $t$.
  \item Total annual cost (€/yr) and LCOE (€/MWh) for the VE scenario.
\end{itemize}

The comparison metric is \textbf{total annualised cost}: capex (annualised),
fixed O\&M, variable O\&M on dispatched generation, net grid expenditure
(purchases minus export revenue), grid running tariffs, and annualised
tilslutningsbidrag. The technology with the lower total cost wins for a given
$(x, \text{year})$ pair. The model is run across years 2022--2025 and
$x \in \{0.25, 0.50, 0.75\}$ to assess robustness to the on-site requirement
and input price variation.

\end{document}
